% auto-generated figure: seq-injury-sub (source: wf)
\begin{figure}[H]
\centering
{\scriptsize
\begin{sequencediagram}
  \newthread{op}{:Operator}
  \newinst[1]{fe}{:Frontend}
  \newinst[1]{gw}{:Gateway}
  \newinst[1]{tm}{:training-match}
  \newinst[1]{mf}{:medical-fitness}
  \newinst[1]{db}{:DB}
  \begin{call}{op}{record INJURY event(player, minute)}{fe}{substitution required}
    \begin{call}{fe}{POST /match-events}{gw}{201}
      \begin{call}{gw}{create(event\_type=INJURY)}{tm}{event}
        \begin{call}{tm}{persist match\_event}{db}{ok}
        \end{call}
        \begin{call}{tm}{POST /injuries(player, status=INJURED)}{mf}{injury}
          \begin{call}{mf}{persist injury (player flipped injured)}{db}{ok}
          \end{call}
        \end{call}
      \end{call}
    \end{call}
  \end{call}
  \begin{call}{op}{record SUBSTITUTION(out=injured, in=sub, minute)}{fe}{updated}
    \begin{call}{fe}{POST /match-events + PUT /match-lineups}{gw}{200}
      \begin{call}{gw}{updateLineup(playerOut@minute, subIn)}{tm}{lineup}
        \begin{call}{tm}{persist match\_event + match\_lineups}{db}{ok}
        \end{call}
      \end{call}
    \end{call}
  \end{call}
\end{sequencediagram}
}
\caption{Sequence diagram for an in-match injury and the forced substitution it triggers, showing how training-match persists the INJURY match event while delegating to medical-fitness to create the real injury record before the operator records the substitution.}
\label{fig-seq-injury-sub}
\end{figure}
